% Methods: what is designed and built so far. Numbers in the testbed tables come from the repository's runs
% (README "Results"); rerun before submission.

\section{Methods}
\label{sec:methods}

\subsection{Problem setting}
A robot arm holds a hand tool and applies it to a workpiece: one arm, one tool, one target. We study three tools
whose loads on the grasp span the regimes a general tool controller must handle:
\begin{itemize}
  \item \textbf{hammer} driving a pre-started nail: an impulse of $\sim$\SI{4}{\milli\second} and
    \SIrange{500}{700}{\newton} at a predictable time;
  \item \textbf{hand saw} crosscutting a board: a periodic load that reverses with every stroke, with binding as an
    untimed event;
  \item \textbf{cordless driver} seating a wood screw or drilling through a board: a steady reaction torque with
    untimed jerks (cam-out, clutch slip, the bit catching, breakthrough).
\end{itemize}
The task goal (nail flush, cut through, screw seated, hole through) must be reached without losing the tool,
without the tool slipping in the grasp, and within the robot's hardware limits.

\subsection{Architecture: anticipate, then react}
The controller has four layers running at different rates (Fig.~\todo{flow chart}).
\begin{description}
  \item[L3, vision-language-action planner (\SIrange{1}{10}{\hertz}).] From camera images, proprioception and a
    command (``drive it flush'', ``tap it in''), it outputs the task phase, the target, and a \emph{prediction} of
    the coming load: the expected grip wrench over the task cycle $\hat{\vect{w}}(\phi)$, the contact time
    $\hat t_c$, the tool mass and an estimate of the grip friction $\hat\mu$.
  \item[L2, learned reactive layer (\SI{1}{\kilo\hertz}).] A small causal network conditioned on the L3 prediction.
    Every millisecond it reads the last \SIrange{16}{32}{\milli\second} of wrist force/torque, pad accelerometers
    and pressure arrays, and outputs corrections to the L1 set points: equilibrium pose, directional stiffness and
    damping, feedforward wrench and grip force, plus estimates of the contact phase and of slip.
  \item[L1, hand-designed controller (\SI{1}{\kilo\hertz}).] Cartesian impedance control with reference
    limiting, impact detection, feedback gating and the hardware limits (Sec.~\ref{sec:l1}). It bounds whatever
    L2 commands.
  \item[L0, tool--target estimator.] Per cycle (strike, stroke, revolution) it estimates task progress (nail depth,
    cut depth, screw seating), the workpiece's resistance and the slip of the tool in the grasp.
\end{description}
The split follows one physical constraint. A hammer blow lasts \SI{4}{\milli\second}, while grip actuators need
\SIrange{50}{185}{\milli\second} to change force, so the grasp cannot be rescued during a blow: it must be set
\emph{before} contact from a prediction (anticipation). Rebound, ringing, slip and jerks unfold over the following
\SIrange{5}{60}{\milli\second}, where a \SI{1}{\kilo\hertz} loop can still act (reaction).

\subsection{L1: impedance control around impacts}
\label{sec:l1}
The arm runs joint-torque control at \SI{1}{\kilo\hertz}:
\begin{equation}
  \vect\tau = J^\top\!\left(K\vect e + D\dot{\vect e} + \vect F_{\mathrm{ff}} + \Lambda\ddot{\vect x}_{\mathrm{ff}}\right)
  + N\vect\tau_{\mathrm{null}} + \vect g(\vect q) + \vect c(\vect q,\dot{\vect q}),
\end{equation}
with $\vect e$ the pose error to the equilibrium (clipped to $\pm\SI{5}{\centi\metre}$, $\pm\SI{0.3}{\radian}$ so
tracking error cannot inject force), $K$ and $D$ diagonal in a task frame, $D=2\zeta\sqrt{K\Lambda}$ by default,
$\Lambda$ the task-space inertia and $N$ the dynamically consistent null-space projector. Stiffness changes are
slew-limited (\SI{5e4}{\newton\per\metre\per\second}) to keep the controller passive. The commanded torque is
clipped to the joint ratings and rate-limited to \SI{1000}{\newton\metre\per\second} (libfranka's limit).

\paragraph{Impacts} Contact is flagged by the first of three detectors: a generalized-momentum observer
\cite{haddadin2017collisions} (gain \SI{400}{\per\second}), the wrist force along the strike axis departing
\SI{35}{\newton} from its \SI{10}{\milli\second} baseline, and the pad accelerometer departing
\SI{60}{\metre\per\second\squared} from its \SI{5}{\milli\second} baseline. The swing reference follows reference
spreading \cite{vansteen2024rs}: an ante-impact reference continued through the expected contact, an interim mode
around the predicted contact time and a post-impact reference started at the flag. For \SI{50}{\milli\second} after
the flag, joint-velocity feedback is replaced by the reference velocity, since ringing makes measured velocity
unreliable \cite{vansteen2024simtoreal,yang2021impactinvariant}.

\paragraph{Swing under hardware limits} The strike speed is capped at a margin (0.8) of the fastest speed the joint
velocity limits allow along the strike path at the strike pose, $v_{\mathrm{cap}}=\min_i
\dot q_i^{\max}/|(J^{+}\vect t)_i|$, and the swing acceleration at 0.6 of what the joint torque limits can push
there. The swing peaks just before the nail and arrives braking, so the torque reversal is under way when the blow
lands; a swing still accelerating at contact keeps driving the arm into the nail for the
$\sim$\SI{170}{\milli\second} a torque-rate-limited joint needs to reverse. The path may be a straight line with the
tool's orientation held, or an arc: the tool turns about a pivot a distance $L$ behind the face along the handle,
so the face moves on a circle and meets the nail square while the shoulder and elbow, not the wrist, drive the
swing.

\subsection{Grip and impedance from a predicted load}
\label{sec:grip}
The same law sets grip force and arm stiffness for all three tools. Each task is a cycle with phase $\phi$
(strike, stroke, revolution). L3 predicts the wrench the tool will put on the grasp over the cycle,
$\hat{\vect w}(\phi)$, split into the force tangential to the pads $\hat f_t$ and the torque about the pad normal
and the handle axis $\hat\tau$. The grip set point is commanded one actuator latency $\tau_a$ ahead:
\begin{equation}
  \begin{aligned}
  F_g(t) = (1+m_k)\,\max\!\Big(&\frac{\hat f_t(\phi(t+\tau_a))}{2\hat\mu},\;
                                  \frac{\hat\tau(\phi(t+\tau_a))}{2\hat\mu r}\Big)\\
         &+ F_{\min} + \Delta F_{\mathrm{L2}}(t),
  \end{aligned}
  \label{eq:grip}
\end{equation}
with $r$ the handle radius, $m_k$ a safety margin, $F_{\min}$ a holding floor and $\Delta F_{\mathrm{L2}}$ the
reactive layer's correction. Loads are routed by how predictable their timing is and how fast they are:
\begin{itemize}
  \item slow, predictable loads (swing inertia, saw strokes, steady drill torque): feedforward from
    $\hat{\vect w}(\phi)$ plus feedback;
  \item fast loads at a predictable time (a blow, a screw head seating): grip and stiffness raised before
    $\hat t_c$; velocity feedback gated after the flag;
  \item fast loads at unpredictable times (cam-out, binding, a bent nail, breakthrough): a standing margin set by
    the prediction's uncertainty, the L2 reaction, and passive pad and wrist compliance for what no actuator can
    follow.
\end{itemize}
For the hammer the grip ramp starts at least \SI{150}{\milli\second} before the predicted contact and peaks
\SI{60}{\milli\second} after it, following human grip during collisions, where grip force rises for
$\sim$\SI{200}{\milli\second} before impact and peaks $\sim$\SI{60}{\milli\second} after peak load
\cite{johansson1988programmed,turrell1999grip,white2011grip}; in humans the same predictive grip--load coupling is
seen in cyclic movements \cite{flanagan1993modulation}. After each cycle $k$ the margin is updated from the measured
slip $s_k$:
\begin{equation}
  m_{k+1} = \operatorname{clip}\!\left(m_k + \gamma\,(s_k - s^\star),\; m_{\min},\; m_{\max}\right),
\end{equation}
and the residual between predicted and measured grip wrench updates $\hat{\vect w}(\phi)$, an iterative-learning
update indexed by phase. Slip is never rescued within a blow; it changes the next cycle.

\subsection{Sensing}
The wrist carries a 6-axis force/torque sensor at \SI{4}{\kilo\hertz}. Each pad of the parallel gripper carries an
$8\times8$ pressure array (\SI{17}{\milli\metre}) at \SI{1}{\kilo\hertz}, with \SI{0.1}{\newton} noise and a
\SI{20}{\newton} per-taxel range, and a 3-axis accelerometer at \SI{8}{\kilo\hertz} ($\pm16\,g$). Joint encoders and
torques run at \SI{1}{\kilo\hertz}. Each sensor model applies a bandwidth filter, decimation, latency, bias, noise,
quantization and saturation.

The dexterous hand wears TaxelScan skins: a 128-taxel sheet on the palm and 32-taxel patches on each finger's
fingertip pad and middle segment, 448 taxels in all, each patch scanned at \SI{1}{\kilo\hertz} by its own RP2350
board. Every taxel is ray-cast onto the palmar surface of its link's detailed mesh, so the patches follow the
fingers' curvature. A contact's normal force is distributed over the taxels as on an elastic foundation (load in
proportion to the skin's compression, found by rays from the taxels to the touching tool), so a handle across the
palm loads a band. The readout models the board: a piezoresistive divider, the 12-bit SAR ADC at
\SI{500}{ksps} (ENOB 9.2), a sequential scan whose per-taxel time skew is kept, calibration residuals, and USB
latency.

\subsection{Simulation testbed}
\label{sec:testbed}
The testbed is built on MuJoCo \cite{todorov2012mujoco} at \SI{8}{\kilo\hertz} (\SI{0.125}{\milli\second},
implicit-fast integration, elliptic friction cones and a no-slip pass, needed for a static grasp to hold). All
contacts are explicit pairs. Every robot and hand comes from its maker's published model at a pinned commit, and its
limits are enforced or monitored as the hardware would:
\begin{itemize}
  \item \textbf{Franka FR3} with torque control: joint torques (\SIrange{12}{87}{\newton\metre}) and the
    \SI{1000}{\newton\metre\per\second} torque rate enforced; joint velocities, ranges and the \SI{3}{\kilo\gram}
    payload monitored;
  \item \textbf{Franka Hand} with compliant \SI{17}{\milli\metre} pads and a grip force loop at \SI{500}{\hertz};
  \item \textbf{WUJI Hand 2} (20 joints, per-joint torque ratings of \SIrange{0.2}{2}{\newton\metre}) holding the
    hammer in a power wrap synthesized offline; joint impedance at \SI{1}{\kilo\hertz}; backdrivable or
    self-locking drives (a ratchet on the joint limits), with hard-stop and gearbox loads monitored against the
    ratings;
  \item \textbf{Dexmate Vega U} humanoid (two 7-joint arms on a fixed pedestal) with a WUJI hand on each arm. As
    on the hardware, its arms take only joint position targets at \SI{100}{\hertz}, tracked by torque-limited PD
    servos; a host loop detects impacts at \SI{1}{\kilo\hertz} and sends differential-IK targets with gravity-droop
    compensation, rate-limited to 90\,\% of the joint velocity limits.
\end{itemize}
Ground truth the sensors never see is logged every step: tool--target contact force, task progress, and the tool's
pose in the grasp (slip).

\paragraph{Plants} MuJoCo cannot remove material, so each workpiece is a small state model whose forces act on the
tool.
\begin{itemize}
  \item \emph{Nail.} A slide joint whose Coulomb resistance grows with depth (\SI{450}{\newton} plus
    \SI{e4}{\newton\per\metre}), with viscous crushing; the hammer--nail contact damping is set from the approach
    speed just before contact (velocity-dependent restitution) and frozen through the blow.
  \item \emph{Saw.} The kerf depth $c$ is a state. At the ends of the stretch of teeth over the board, the kerf
    bottom supports the blade with $k_n\,(e-c)^+$ plus damping ($e$: the tooth's depth below the surface); the
    teeth resist the stroke with $\mu_c F_n$ plus drag; on the cutting stroke the kerf deepens at
    $\dot c = F_n|v_s|/k_c(c)$, with a knot raising $k_c$ over a depth range; once the kerf is started its walls push
    the blade back beyond the side clearance, adding binding friction.
  \item \emph{Driver and screw.} A motor $\tau_m = \tau_s(u - \omega/\omega_0)$ with a soft-start trigger and an
    electronic brake drives the spindle through a clutch that slips above its setting with a detent ripple. The
    screw advances one pitch per turn against a torque growing with depth and rising steeply once the head is flush.
    The bit cams out of the recess when the torque exceeds $c_{\mathrm{cam}}F_a(1-\theta/\theta_{\max})$ ($F_a$ the
    axial push, $\theta$ the bit's misalignment), with a kick and ripple; repeated cam-outs strip the recess. The
    housing carries the output torque plus the rotor's acceleration.
  \item \emph{Driver and hole.} Thrust above a threshold sets the feed per revolution and the torque; over the last
    \SI{2}{\milli\metre} the support fades and the bit grabs; once through, the tool loses its support.
\end{itemize}

\paragraph{Tasks with scripted layers} Until the learned layers exist, scripted stand-ins issue the same L2 commands.
The hammer swing plans a limit-aware strike, ramps the grip before contact and raises the margin after a slip. The
saw lowers the teeth onto the board and strokes at \SI{1.2}{\hertz} ($\pm\SI{7}{\centi\metre}$) with a
\SI{25}{\newton} push and a soft spring along the cut. The driver brings the bit to the screw, pushes
\SI{70}{\newton} along the bit (soft along it, stiff across it), pulls the trigger once the wrist force reads the
push, and stops when the tool's own telemetry reports the screw seated or the bit through.

\paragraph{Replaying recorded motions} No public dataset records impact or grip forces in tool use, so recorded
motions supply kinematics and the testbed supplies the contact. Human hammering from the Adroit/DAPG demonstrations
(VR and data glove, \SI{100}{\hertz}) and tool trajectories tracked from human videos in DexToolBench
(\SI{3}{\hertz}) are retargeted face-first, so every recorded contact becomes a square blow on our nail head: the
recorded strike is placed on our nail along our strike axis, each contact is aimed at the head, the motion across the
nail axis is blended out within \SI{80}{\milli\second} of a contact, the face is turned square at every contact, a
soft clamp keeps the face \SI{3}{\milli\metre} short of the head between blows, and the grasp pose follows from our
hammer's face offset. A time warp slows the path only where it exceeds the arm's speed or acceleration caps (ignoring
the deceleration within \SI{40}{\milli\second} of a recorded contact, which the nail produces), and damped
least-squares IK along the path checks reach and joint speeds. On the FR3 the 22 Adroit demonstrations that contain
blows replay within all limits and track to \SIrange{2.5}{4.4}{\milli\metre} RMS; 79 blows land for 75 recorded
contacts, one off-centre and none against the shank. But they need 2.2--4.6$\times$ the recorded time: a human hand
accelerates far beyond what the arm can, so the replayed blows land at \SIrange{0.07}{0.61}{\metre\per\second}. On
the Vega U, 21 of the 22 replay (73 blows for 72 contacts, \SIrange{155}{370}{\newton}); its position interface lags,
so 18 blows end with the tool resting on the nail for a moment.

\subsection{Preliminary testbed results}
Table~\ref{tab:hammer} summarizes ten-strike hammer episodes at \SI{8}{\kilo\hertz}, and Table~\ref{tab:tools}
the saw and driver tasks. They come from the scripted layers and fixed grip settings; they establish the testbed and
the baselines the learned layers will be compared against.

\begin{table}[t]
  \centering
  \caption{Hammer, ten strikes (\SI{8}{\kilo\hertz}, seed 0, scripted layers).}
  \label{tab:hammer}
  \resizebox{\columnwidth}{!}{%
  \begin{tabular}{@{}lcccc@{}}
    \toprule
    & FR3 + & FR3 + & FR3 + WUJI & Vega U + \\
    & Franka Hand & WUJI & self-locking & WUJI \\
    \midrule
    Strikes on nail    & 10/10 & 10/10 & 10/10 & 10/10 \\
    Nail driven (mm)   & 13.8 & 12.3 & 15.6 & 10.2 \\
    Face speed (m/s)   & 1.3--1.8 & 1.2--2.2 & 1.2--1.8 & 1.08--1.13 \\
    Flag latency (ms)  & 0.4--1.1 & 0.25--1.1 & 0.25--1.1 & 0.6--1.75 \\
    Tilt/strike (deg)  & 0.5--2.9 & 1.9--30.7 & 0.1--0.5 & 0.3--2.8 \\
    Arm limits         & held & J6 speed 1.14$\times$ & J6 speed 1.31$\times$ & held \\
    Hand joint loads   & --- & stops 16$\times$ & gearbox 32$\times$ & stops 4$\times$ \\
    \bottomrule
  \end{tabular}}
\end{table}

\begin{table}[t]
  \centering
  \caption{Saw and driver on the FR3 + Franka Hand (\SI{8}{\kilo\hertz}, scripted layers, grip \SI{60}{\newton}
  per pad).}
  \label{tab:tools}
  \resizebox{\columnwidth}{!}{%
  \begin{tabular}{@{}lccc@{}}
    \toprule
    & Saw, crosscut & Driver, screw & Driver, hole \\
    \midrule
    Goal                  & \SI{20}{\milli\metre} cut & seated & \SI{38}{\milli\metre} through \\
    Reached at (s)        & 23.1 (54 strokes) & 2.16 & 7.14 \\
    Peak load             & \SI{49}{\newton} normal &  \SI{1.88}{\newton\metre}, \SI{87}{\newton} & \SI{2.41}{\newton\metre}, \SI{90}{\newton} \\
    Net slip in grasp     & \SI{1.0}{\milli\metre}, \SI{4.8}{\degree} & \SI{0.2}{\milli\metre}, \SI{2.4}{\degree} & \SI{0.1}{\milli\metre}, \SI{0.6}{\degree} \\
    Untimed events        & binding (\SI{23}{\newton}) & 0 cam-outs & breakthrough \\
    Arm limits            & held & held & held \\
    \bottomrule
  \end{tabular}}
\end{table}

Three findings shape the method. (i) On the FR3, enforcing the torque-rate limit made an accelerate-through-contact
swing infeasible (it demanded up to 43$\times$ the limit), which motivated the braking, limit-capped swing. (ii) A
dexterous power wrap at the WUJI hand's rated torques does not hold a blow: the handle turns
\SIrange{2}{31}{\degree} per strike and finger joints hit their stops at up to 16$\times$ their rating; with
self-locking drives the tool holds (\SIrange{0.1}{0.5}{\degree}) but the gearboxes carry up to 32--40$\times$ their
rating. Holding by squeeze alone is therefore not available to a dexterous hand; holding must come from geometry,
timing and locking. (iii) At a fixed grip, a parallel gripper holding a driver by its body yields \SI{2.4}{\degree} about the bit when
the clutch ratchets at seating, and a saw handle flexes the pads by up to \SI{9}{\degree} each stroke and creeps
\SI{4.8}{\degree} over 27 strokes: the jerk and periodic loads that Eq.~\eqref{eq:grip} and the cycle-to-cycle
margin update are meant to anticipate.

\subsection{Planned evaluation}
\todo{metrics: progress per cycle, slip per cycle, grip effort, reaction latency from contact to actuator; baselines:
fixed high grip at equal mean force, reactive-only slip control, hand-designed L1 only, a slow (24--30\,Hz) fast
layer as in \cite{xue2025rdp}, a grasp optimizer as in \cite{gupta2026grasptoact,trupin2025physics}; ablations:
anticipation lead time, L2 rate (30, 100, 300, 1000\,Hz), sensing rate and modality, VLA prediction removed;
sim-to-real: a force-truth rig (force plate, nail-depth encoder, hammer IMU) to check the simulated impulses.}
